
\chapter[Na\"{i}ve Bayes (1 of 3)]
{\LARGE ML classifiers: Na\"{i}ve Bayes (1 of 3)}
\label{ch:naiveBayes1}

\index{machine learning (ML)}

For some folks, the coolest and most exciting parts of Data Science are the
\textbf{machine learning (ML)} algorithms. These are the under-the-hood
techniques for doing prediction and analytics. We'll look at a few of them
briefly in this class, and you'll delve into them more fully in later courses.

\index{classifier}
\index{label}
\index{class (ML classifiers)}

Recall that a \textbf{classifier} is an algorithm/program to predict which of
several \textbf{labels} (or \textbf{classes}) to assign to a data point.
Perhaps we're trying to predict what political party someone is affiliated
with, and our labels are ``Democrat,'' ``Republican,'' and ``Green.'' Perhaps
we're trying to predict whether someone will default on a loan, and the labels
are ``safe'' and ``high-risk.''

\index{labeled example}
\index{unlabeled example}
\index{example, labeled vs.~unlabeled}
\index{training data}

There are many different algorithms for classification, but they all use a set
of \textbf{labeled examples} as ``\textbf{training data}.'' The idea is that by
scrutinizing a bunch of examples for which we \textit{know} the label, we can
intelligently figure out what label to assign to new, \textbf{unlabeled
examples}.

\index{feature}
\index{target}

The attributes of our training data (and new data) that we'll use for
prediction are sometimes called \textbf{features}. If our training data is in
a DataFrame, the features are essentially all the relevant columns except for
one, which is the \textbf{target}. Our goal, given a new data point that has
values for all the features, is to label its target with the label it's most
likely to have, assuming the training data are representative examples. 

\index{performance (of a classifier)}
\index{test data}

A classifier's \textbf{performance} can be measured in several ways, but the
most straightforward is to score how many unlabeled examples it gives the
correct label to. This of course begs the question: if the example is
\textit{un}labeled, how can we know whether or not the classifier's label was
correct? What we have to do in practice is set aside some of our (precious)
labeled data as a set of \textbf{test data}. We sort of pretend we don't know
the answers for the test data, feed them to the algorithm which has been
trained on the \textit{training} points, and then score how well it did.

\index{Na\"{i}ve Bayes classifier}
\index{baseline}

One of the simplest (and quickest) classifiers around is called the
\textbf{Na\"{i}ve Bayes classifier}, for reasons which will become apparent. In
some ways, it's the rock-bottom, no-frills default algorithm for
classification. One reason it's important is that it's simple and actually
works surprisingly (even embarrassingly) well in many cases. Another reason is
that even when it doesn't turn out to have the best performance, it's still a
\textbf{baseline} which other algorithms' performance can be compared to. A
researcher might report: ``our whiz-bang algorithm X got 89.4\% of examples
correct on this test set, compared with 83.7\% for straightforward Na\"{i}ve
Bayes.''

\section{The Na\"{i}ve Bayes algorithm}

\index{spycam}
\index{security camera}

Suppose we're designing a security system for a Galactic Cruiser. One thing we
need to do is identify possible evildoers on our security cameras. Given what
we can observe about them, can we tell good guys from bad guys?

\index{training data}
\index{security\_cam@\texttt{security\_cam.csv}}
\label{jediTraining}

\begin{Verbatim}[fontsize=\small,samepage=true,frame=leftline,framesep=5mm,framerule=1mm]
     face  height  demeanor lightsaber  type
    clean   short  menacing      green  Sith
    beard  medium      calm      green  Jedi
    clean    tall  menacing        red  Sith
    beard  medium      calm       blue  Jedi
    beard  medium      calm       blue  Jedi
    clean  medium      calm       blue  Sith
    clean    tall  menacing        red  Jedi
    clean  medium  menacing        red  Sith
    clean   short      calm       blue  Jedi
                    . . .
\end{Verbatim}

\index{feature}
\index{label}
\index{target}
\index{Jedi}
\index{Sith}

The training data above, in a \texttt{DataFrame} called ``\texttt{t}'' (for
``training''), has four features: whether the \texttt{face} of the person in
question has a \texttt{beard} or is \texttt{clean}-shaven, what their
\texttt{height} is (clustered in three broad ranges), whether their
\texttt{demeanor} is \texttt{menacing} or \texttt{calm}, and what color
\texttt{lightsaber} they wield. The final column is our label: this person is
known to be a \texttt{Jedi} (good guy) or a \texttt{Sith} (bad guy).

A test point that we might want to evaluate looks like this:
\begin{Verbatim}[fontsize=\small,samepage=true,frame=leftline,framesep=5mm,framerule=1mm]
     face  height  demeanor lightsaber  type
    beard   short  menacing        red   ?? 
\end{Verbatim}


In other words, a short, menacing, bearded person with a red lightsaber has
just showed up on security camera \#305. Should we sound the alarm? That, of
course, depends upon their (as yet unknown) \texttt{type}. We want to predict
this based on our past examples.

\index{conditional probability}
\bigskip

The \textbf{Na\"{i}ve Bayes} classification algorithm applies conditional
probability to this question. Really, what we're asking is: ``which of the
following two quantities is greater:

\begin{compactenum}
\item P(Jedi\ $|$\ beard,short,menacing,red),  or
\item P(Sith\ $|$\ beard,short,menacing,red)?''
\end{compactenum}
 
Note that we're combining joint and conditional probabilities here. Given that
we have the \textit{joint} event ``bearded, short, menacing, and red,'' what's
the \textit{conditional} probability of them being a Jedi (or Sith)? If we
estimate the first as .8 and the second as .2, for instance, we won't sound the
alarm.

\section{Using Bayes' rule}

\index{Bayes' Rule (probability)}

With Bayes' rule, we can rewrite these as:

\footnotesize
\begin{center}
(1) Estimated Jedi probability:

\boxed{\textrm{P(Jedi | beard,sh,menace,red)} =
\frac{\textrm{P(beard,sh,menace,red | Jedi)} \cdot \textrm{P(Jedi)}}{\textrm{P(beard,sh,menace,red)}}}
\end{center}
\normalsize

and
\footnotesize
\begin{center}
(2) Estimated Sith probability:

\boxed{\textrm{P(Sith | beard,short,menace,red)} =
\frac{\textrm{P(beard,short,menace,red | Sith)} \cdot \textrm{P(Sith)}}{\textrm{P(beard,short,menace,red)}}}
\end{center}
\normalsize

Note that the two equations are identical except for the words ``Jedi'' and
``Sith,'' which are interchanged. Let's focus on equation (1) first.

Bayes' Rule has allowed us to turn things on their head. In essence, instead of
asking ``if they have these physical features, would they likely be a Jedi?''
we're asking ``if they're a Jedi, would they be likely to have these physical
features?'' 

\index{prior (probability)}

\index{numerator}

The right-hand side of equation (1) has three parts: two in the numerator, and
one in the denominator. Ignore the denominator for a moment. The second part of
the numerator, P(Jedi), is our \textbf{prior} probability. That's easy to
estimate from the training data: we just count up what percentage of our
samples are Jedi overall.

The first half of the numerator needs more work. It reads:

\vspace{-.2in}
\begin{align*}
P(\textrm{beard},\textrm{short},\textrm{menacing},\textrm{red}\ |\ \textrm{Jedi}) 
\end{align*}

\index{independence (probability)}

Now \textbf{\textit{\underline{if}}} \ it just so happened that those four
features were \textit{independent} of one another, this would actually be
super easy to compute. We could just do this:

\vspace{-.2in}
\begin{align*}
P(\textrm{beard} \ |\ \textrm{Jedi}) \cdot
P(\textrm{short} \ |\ \textrm{Jedi}) \cdot
P(\textrm{menacing} \ |\ \textrm{Jedi}) \cdot
P(\textrm{red} \ |\ \textrm{Jedi})
\end{align*}

Under this outrageous assumption, we've split the expression into \textit{each
feature individually.} This is called \textbf{the Na\"{i}ve Bayes assumption}.
Of course it's probably \textit{not} true that beardedness, height, demeanor,
and lightsaber color are all independent of each other; for one thing, if some
or all of these features are indicative of Jedi-ness (which we hope), then they
certainly \textit{will} vary together! But the Na\"{i}ve Bayes classifier says
``aw heck, let's just throw caution to the wind and see how things work out.''

All these components are now trivially obtainable from the training data. For
example, what's a good estimate of P(beard $|$ Jedi)? Simple: look at
\textit{only the Jedi} in the training data, and count up how many
``\texttt{beard}s'' and ``\texttt{clean}s'' there are. The fraction that are
\texttt{beard}s is our estimate of P(beard $|$ Jedi).\footnote{If you have
trouble seeing this, go back to the definition of conditional probability on
p.~\pageref{condProbDef}. We're really calculating
$\dfrac{P(\textrm{beard},\textrm{Jedi})}{P(\textrm{Jedi})}$ here, which is
``the number of bearded Jedi divided by the number of total Jedi.''}

To make it concrete, in our training data on p.~\pageref{jediTraining}, there
are a total of six Jedi (count them!), of whom four have beards (count them!).
Therefore, our estimate for P(beard $|$ Jedi) is $\frac{4}{6} = \frac{2}{3}$.
In other words, based on this training data, we guesstimate that if you're a
Jedi, it's about 67\% likely that you're bearded.

\section{A numerical example}

\label{jediNumerical}
\index{numerator}

Let's work out the numerical answer (numerator only) for the training data
listed earlier. First, considering \textit{only} the five \texttt{Jedi} rows
gives us our Jedi numerator:

\smallskip
\footnotesize
% Whew! See array package documentation.
\hspace{-.2in}
\begin{tabular}{>{$}c<{$}@{$\ \cdot\ $}>{$}c<{$}@{$\ \cdot\ $}>{$}c<{$}@{$\
\cdot\ $}>{$}c<{$}@{$\ \cdot\ $}>{$}c<{$}>{$}l<{$}}
P(\textrm{beard}|\textrm{Jedi}) &
P(\textrm{short}|\textrm{Jedi}) &
P(\textrm{menace}|\textrm{Jedi}) &
P(\textrm{red}|\textrm{Jedi}) & P(\textrm{Jedi}) &= \\
\frac{4}{6} & \frac{1}{6} & \frac{2}{6} & \frac{1}{6} &
\frac{6}{11} &= .0034
\end{tabular}
\normalsize

\bigskip
Now let's do the same thing for the five \texttt{Sith} rows to get the Sith
numerator:

\footnotesize
\smallskip
\hspace{-.2in}
\begin{tabular}{>{$}c<{$}@{$\ \cdot\ $}>{$}c<{$}@{$\ \cdot\ $}>{$}c<{$}@{$\
\cdot\ $}>{$}c<{$}@{$\ \cdot\ $}>{$}c<{$}>{$}l<{$}}
P(\textrm{beard}|\textrm{Sith}) &
P(\textrm{short}|\textrm{Sith}) &
P(\textrm{menace}|\textrm{Sith}) &
P(\textrm{red}|\textrm{Sith}) & P(\textrm{Sith}) &= \\
\frac{1}{5} & \frac{1}{5} & \frac{4}{5} & \frac{2}{5} &
\frac{5}{11} &= .0058\\
\end{tabular}
\normalsize

\bigskip
Look very very carefully at the training data on p.~\pageref{jediTraining} and
confirm that you agree with the above numbers. Using a pencil to write marks
on the training data often helps you to keep everything straight.

\subsection{Throw away the denominator}

\index{denominator}
\index{throwing away the denominator}

Okay, those are the two numerators. To get true probabilities, we have to
divide by the denominators, of course. But the nice thing is: \textit{the
denominators are the \underline{same} in the two cases!} For both expressions,
we would need to divide by P(\textrm{beard}, \textrm{short}, \textrm{menacing},
\textrm{red}) (in words, ``what's the probability that any random person would
be a bearded, short, menacing person with a red lightsaber?'') This of course
has the same value even if we compute it twice. For this reason, we
\textit{throw away the denominator} since in the end all we really care about
is ``which is higher, P(Jedi $|$ \textsl{stuff}) or P(Sith $|$
\textsl{stuff})?''

In this case, P(Sith $|$ \textsl{stuff}) is the higher of the two quantities
(.0058 vs.~.0034) and so we boldly predict: ``this person looks like a Sith!
Sound the alarm.''

\bigskip
All that counting and dividing, of course, is a pain. So we'll get Python to
automate it for us. That's the subject of the next chapter.
